\begin{lemma}
For any prefix \(p\) and child \(x\), put
\[
  r=r^p(x.2),\qquad \ell=\ell_r^p.
\]
Then \(U_\ell^p\) is nonempty; in Lean's real-valued cardinal notation,
\[
  1\le |U_\ell^p|.
\]
\end{lemma}
\begin{proof}
The point \(x.2\) lies in \(U_r^p\), because its rank is exactly \(r\).  Thus
\(U_r^p\) has positive cardinality.  Applying
\(\noderef{StarProductLayerChoicePositive}\) shows that the selected layer
\(Z_\ell^p\) has positive cardinality.  Since \(Z_\ell^p\subseteq U_\ell^p\),
the set \(U_\ell^p\) also has positive cardinality, equivalently real
cardinality at least \(1\).
\end{proof}
